% ==============================================================================
% taxonomy.tex
% ==============================================================================

\section{Taxonomy of Methods}
\label{sec:taxonomy}

We organise the surveyed methods into four major architectural families, reflecting both the historical evolution and the distinct design philosophies in the field.

\begin{enumerate}
  \item \textbf{GAN-Based Methods} (Section~\ref{sec:gan_methods}): These approaches employ generative adversarial networks to learn a mapping from 2D X-ray projections to 3D CT volumes. A generator network synthesises the volume, while a discriminator provides adversarial supervision to improve perceptual quality.

  \item \textbf{Transformer-Based Methods} (Section~\ref{sec:transformer_methods}): Leveraging self-attention mechanisms and vision transformers, these methods capture long-range spatial dependencies that are difficult for purely convolutional architectures to model.

  \item \textbf{Neural Field Methods} (Section~\ref{sec:neural_field_methods}): Inspired by neural radiance fields (NeRF), these approaches represent the CT volume as a continuous implicit function parameterised by a neural network, enabling reconstruction at arbitrary resolution.

  \item \textbf{Diffusion-Based Methods} (Section~\ref{sec:diffusion_methods}): The most recent family applies denoising diffusion probabilistic models or score-based generative models to the reconstruction task, treating 3D volume synthesis as a conditional generation problem.
\end{enumerate}

Figure~\ref{fig:taxonomy} (to be added) illustrates the proposed taxonomy and the relationships among these families.

% TODO: Insert taxonomy diagram (figures/taxonomy.pdf or .png).
% \begin{figure}[t]
%   \centering
%   \includegraphics[width=\linewidth]{figures/taxonomy_placeholder}
%   \caption{Proposed taxonomy of few-view X-ray-to-CT reconstruction methods.}
%   \label{fig:taxonomy}
% \end{figure}
